\chapter*{Lời Mở Đầu}
\addcontentsline{toc}{chapter}{Lời Mở Đầu}
\markboth{Lời Mở Đầu}{Lời Mở Đầu}

\begin{truyen}{Lời Mở Đầu}{Opening Words}
\chuhoa{T}{hiên nhiên ngàn đời nay vẫn đang dạy ta — chỉ là ta hay quên lắng nghe.}
\textit{(Nature has been teaching us for thousands of years --- we simply often forget to listen.)}

Quyển ba của bộ sách chọn góc nhìn từ thế giới tự nhiên: những loài vật nhỏ bé, những quy luật vũ trụ, những hiện tượng thường ngày. Tất cả đều có bài học ẩn bên trong.
\textit{(The third volume of the series chooses the perspective of the natural world: small creatures, universal laws, everyday phenomena. Each carries a hidden lesson.)}

Học sinh đọc những chuyện này không chỉ để biết thêm về thiên nhiên, mà để nhìn lại bản thân mình qua tấm gương của muôn loài.
\textit{(Students read these stories not only to learn more about nature, but to look at themselves through the mirror of all living things.)}

Con nhện chăng tơ dạy về kiên nhẫn. Con kiến tha hồ dạy về tích lũy. Nước chảy đá mòn dạy về sự bền bỉ.
\textit{(The spider spinning its web teaches patience. The ant gathering food teaches accumulation. Water wearing away stone teaches perseverance.)}
\end{truyen}

\begin{baihoc}
Những bài học từ thiên nhiên không cần phải học --- chúng chỉ cần được nhìn thấy.
\textit{(Lessons from nature do not need to be studied --- they only need to be seen.)}
\end{baihoc}

\begin{tcolorbox}[colback=truyenblue!6,colframe=truyenblue,coltitle=white,fonttitle=\bfseries,title={Easy English},arc=4pt,boxrule=1pt,breakable]
\textbf{Short English to remember:}

Nature teaches, if we choose to observe.

Small creatures carry big wisdom.
\end{tcolorbox}

\begin{tcolorbox}[colback=truyengold!10,colframe=truyengold!70!black,coltitle=black,fonttitle=\bfseries,title={Tự Học Nhanh},arc=4pt,boxrule=1pt,breakable]
1. Trong thiên nhiên, con vật nào em thấy mình giống nhất? \textit{In nature, which creature do you feel most resembles you?}

2. Quy luật tự nhiên nào em nghĩ có thể áp dụng vào việc học? \textit{Which natural law do you think can be applied to your studies?}

3. Hôm nay em sẽ quan sát một điều gì đó trong thiên nhiên và suy nghĩ về ý nghĩa của nó? \textit{Today, will you observe something in nature and reflect on its meaning?}
\end{tcolorbox}

\ngancach
